% GENERATED FILE. Edit canonical lessons, then run npm run generate:books.
% canonical-source-hash: fnv1a64:295c44dc858a33d3
% canonical-lessons: ZH-C32-mama, ZH-W32-ma-mum, ZH-C32-baba, ZH-C32-gege, ZH-C32-jiejie, ZH-C32-didi

\chapter{Mum and Dad}
\label{ch:zh-mum-and-dad}
\hlchaptermodality{32}

\begin{chapteropening}
\textbf{By the end of this chapter:} \emph{I can name my parents, brothers and sisters.}

Complete the last lesson of chapter 32: I can name my parents, brothers and sisters.
\end{chapteropening}

\section[\texorpdfstring{\zh{妈妈} (māma)}{妈妈 (māma)}]{\zh{妈妈} (māma) — mum; mother}

\label{lesson:ZH-C32-mama}



\begin{quote}
Before the new one: say the Mandarin for we, then the Mandarin for you all.
\end{quote}

\subsection*{You'll want to know: \zh{妈妈}}
\textbf{\zh{妈妈}} — \emph{māma} — \textquotedblleft{}mum; mother\textquotedblright{}.

\begin{quote}
\textbf{\zh{我妈妈很好}} \emph{wǒ māma hěn hǎo} — My mum is well.
\end{quote}

The new character, in pieces: \zh{妈} — \zh{女}, the woman piece, on the left, and on the right a three-stroke piece read mǎ, the word for horse, which gives the sound as it does in \zh{吗}. The next lesson puts it in your hand.

\begin{cousinweb}
\textbf{\zh{妈妈}} (\emph{māma}) is written with a character the Guangyun records as \emph{muX} (Guangyun, in Baxter's transcription), glossed 'mother'. The modern sound is not what \emph{muX} would regularly become, so whether mā is inherited from it, or is a nursery word written with this character, is not settled.
\end{cousinweb}

\begin{grammarlens}[title={What you've built}]
One more word about people.
\end{grammarlens}

\subsection*{Your turn}
\begin{itemize}
\raggedright
  \item \emph{Say it:} \emph{māma}
  \item \emph{Say it:} \emph{māma}, once more
  \item \emph{Say it:} \emph{wǒ māma hěn hǎo} — My mum is well.
  \item \emph{Recall:} say the Mandarin for we, then the Mandarin for you all, then say \emph{māma} again
\end{itemize}

\begin{tcolorbox}[breakable,colback=teal!4,colframe=teal!35!black,title={Before you move on}]
What does \zh{妈妈} mean? (\textquotedblleft{}Mum; mother\textquotedblright{}.) Say it once more.
\end{tcolorbox}
\section[\texorpdfstring{\zh{妈} (mā)}{妈 (mā)}]{\zh{妈} — the woman piece, then a horse}

\label{lesson:ZH-W32-ma-mum}



\begin{quote}
The word first, then the character inside it.

\begin{itemize}
\raggedright
  \item \emph{Recall:} say \emph{māma} — \textbf{R1}, one lesson back, before it decays
\end{itemize}

\begin{quote}
\zh{女} \zh{吗}
\end{quote}

\emph{Read it:} each one aloud, with nothing said first — \emph{nǚ}, \emph{ma}
\end{quote}

\subsection*{Script — the woman piece, then a horse}
\hlblockfigure{\detokenize{figures/ZH-W32-ma-mum-filmstrip.pdf}}{How \zh{妈} is written, stroke by stroke}

\begin{quote}
\zh{妈}
\end{quote}

\textbf{Two pieces, left before right.}

\textbf{The left piece, \zh{女}, three strokes}, as you wrote it alone, but narrow: the fold that drops and turns, the long left-falling stroke, and last the horizontal, which here rises to the right so it does not cross into the next piece.

\textbf{The right piece, the horse, three strokes.} A horizontal that turns down. A vertical that folds right, folds down again and hooks at the foot. Then a rising stroke across the bottom.

3 + 3 = six. The same horse stands on the right of \zh{吗}, which is why the two sound alike.

\subsection*{Your turn}
\begin{itemize}
\raggedright
  \item \emph{Write it:} \zh{妈}, all six, with the model in front of you
  \item \emph{Cover:} write it again from memory, and count the pieces, not the strokes
\end{itemize}

\begin{tcolorbox}[breakable,colback=teal!4,colframe=teal!35!black,title={Before you move on}]
How many pieces in \zh{妈}? (\textbf{Two} — \zh{女}, then the horse.) What does the last stroke of \zh{女} do here? (\textbf{It rises to the right}.) Which character you know has the same right piece? (\textbf{\zh{吗}}.)

Source: \href{https://raw.githubusercontent.com/chanind/hanzi-writer-data/68d10a4b21150cae5e1ebbd223eed289cf32d90c/data/\%E5\%A6\%88.json}{Hanzi Writer Data \zh{妈}, pinned stroke paths and medians}.
\end{tcolorbox}
\section[bàba]{bàba — dad; father}

\label{lesson:ZH-C32-baba}



\begin{quote}
Before the new one: say the Mandarin for you all, then the Mandarin for mum.
\end{quote}

\subsection*{You'll want to know: bàba}
\emph{bàba} — \textquotedblleft{}dad; father\textquotedblright{}.

\begin{quote}
\emph{wǒ bàba xǐhuan kàn shū} — My dad likes reading.
\end{quote}

Its characters come later; for now it is a word for the ear.

\begin{cousinweb}
\emph{bàba} is inherited from Middle Chinese \emph{baX} (Guangyun, in Baxter's transcription), where the character is glossed 'father'. Spoken, it is doubled, like the word for mum.
\end{cousinweb}

\begin{grammarlens}[title={What you've built}]
One more word about people.
\end{grammarlens}

\subsection*{Your turn}
\begin{itemize}
\raggedright
  \item \emph{Say it:} \emph{bàba}
  \item \emph{Say it:} \emph{bàba}, once more
  \item \emph{Say it:} \emph{wǒ bàba xǐhuan kàn shū} — My dad likes reading.
  \item \emph{Recall:} say the Mandarin for you all, then the Mandarin for mum, then say \emph{bàba} again
\end{itemize}

\begin{tcolorbox}[breakable,colback=teal!4,colframe=teal!35!black,title={Before you move on}]
What does bàba mean? (\textquotedblleft{}Dad; father\textquotedblright{}.) Say it once more.
\end{tcolorbox}
\section[gēge]{gēge — elder brother}

\label{lesson:ZH-C32-gege}



\begin{quote}
Before the new one: say the Mandarin for mum, then the Mandarin for dad.
\end{quote}

\subsection*{You'll want to know: gēge}
\emph{gēge} — \textquotedblleft{}elder brother\textquotedblright{}.

\begin{quote}
\emph{wǒ yǒu yí ge gēge} — I've got an older brother.
\end{quote}

Its characters come later; for now it is a word for the ear.

\begin{cousinweb}
\emph{gēge} is inherited from Middle Chinese \emph{ka} (Guangyun, in Baxter's transcription). The dictionary notes that the character was once written for 'song', and that in its own day people used it to call an elder brother.
\end{cousinweb}

\begin{grammarlens}[title={What you've built}]
One more word about people.
\end{grammarlens}

\subsection*{Your turn}
\begin{itemize}
\raggedright
  \item \emph{Say it:} \emph{gēge}
  \item \emph{Say it:} \emph{gēge}, once more
  \item \emph{Say it:} \emph{wǒ yǒu yí ge gēge} — I've got an older brother.
  \item \emph{Recall:} say the Mandarin for mum, then the Mandarin for dad, then say \emph{gēge} again
\end{itemize}

\begin{tcolorbox}[breakable,colback=teal!4,colframe=teal!35!black,title={Before you move on}]
What does gēge mean? (\textquotedblleft{}Elder brother\textquotedblright{}.) Say it once more.
\end{tcolorbox}
\section[jiějie]{jiějie — elder sister}

\label{lesson:ZH-C32-jiejie}



\begin{quote}
Before the new one: say the Mandarin for dad, then the Mandarin for elder brother.
\end{quote}

\subsection*{You'll want to know: jiějie}
\emph{jiějie} — \textquotedblleft{}elder sister\textquotedblright{}.

\begin{quote}
\emph{wǒ jiějie shì dàxuésheng} — My older sister is at university.
\end{quote}

Its characters come later; for now it is a word for the ear.

\begin{cousinweb}
\emph{jiějie} is inherited from Middle Chinese \emph{tsjaeX} (Guangyun, in Baxter's transcription), where it is glossed as the word one people used for mother. How it came to mean elder sister is not settled.
\end{cousinweb}

\begin{grammarlens}[title={What you've built}]
One more word about people.
\end{grammarlens}

\subsection*{Your turn}
\begin{itemize}
\raggedright
  \item \emph{Say it:} \emph{jiějie}
  \item \emph{Say it:} \emph{jiějie}, once more
  \item \emph{Say it:} \emph{wǒ jiějie shì dàxuésheng} — My older sister is at university.
  \item \emph{Recall:} say the Mandarin for dad, then the Mandarin for elder brother, then say \emph{jiějie} again
\end{itemize}

\begin{tcolorbox}[breakable,colback=teal!4,colframe=teal!35!black,title={Before you move on}]
What does jiějie mean? (\textquotedblleft{}Elder sister\textquotedblright{}.) Say it once more.
\end{tcolorbox}
\section[dìdi]{dìdi — younger brother}

\label{lesson:ZH-C32-didi}



\begin{quote}
Before the new one: say the Mandarin for elder brother, then the Mandarin for elder sister.
\end{quote}

\subsection*{You'll want to know: dìdi}
\emph{dìdi} — \textquotedblleft{}younger brother\textquotedblright{}.

\begin{quote}
\emph{wǒ dìdi bā suì} — My little brother is eight.
\end{quote}

Its characters come later; for now it is a word for the ear.

\begin{cousinweb}
\emph{dìdi} is inherited from Middle Chinese \emph{dejX} (Guangyun, in Baxter's transcription), which the dictionary explains as the brother born later.
\end{cousinweb}

\begin{grammarlens}[title={What you've built}]
One more word about people.
\end{grammarlens}

\subsection*{Your turn}
\begin{itemize}
\raggedright
  \item \emph{Say it:} \emph{dìdi}
  \item \emph{Say it:} \emph{dìdi}, once more
  \item \emph{Say it:} \emph{wǒ dìdi bā suì} — My little brother is eight.
  \item \emph{Recall:} say the Mandarin for elder brother, then the Mandarin for elder sister, then say \emph{dìdi} again
\end{itemize}

\begin{tcolorbox}[breakable,colback=teal!4,colframe=teal!35!black,title={Before you move on}]
What does dìdi mean? (\textquotedblleft{}Younger brother\textquotedblright{}.) Say it once more.
\end{tcolorbox}
